\section{Introduction}
\label{sec:introduction}

Let \(\SPD^n\) denote the cone of real symmetric positive definite matrices.
A full-column-rank matrix \(A\in\R^{n\times m}\) defines the positive
compression
\[
  \pi_A:\SPD^n\longrightarrow\SPD^m,
  \qquad
  \pi_A(P)=A^\top P A.
\]
Given a reference \(P_0\in\SPD^n\) and a target \(R\in\SPD^m\), the basic
matrix problem studied here is
\begin{equation}
  \min_{P\in\SPD^n}
  \bigl\{\dai(P_0,P):A^\top P A=R\bigr\}.
  \label{eq:intro-nearest-completion}
\end{equation}
The constraint specifies only the compression of \(P\), so its feasible set
is a noncompact positive-definite fiber.  We ask when this fiber has a unique
affine-invariant nearest point, whether the constrained value reduces exactly
to an \(m\)-dimensional problem, and how the selected completion changes when
the matrix channel itself moves.

This compression problem is also the matrix core of partial-information
training geometry.  Adaptive optimization is often described as a choice of
positive geometry: natural gradient uses the Fisher metric
\citep{amari1998natural}; K-FAC and related methods use structured Fisher or
Gauss--Newton approximations \citep{martens2015kfac,george2018ekfac}; AdaGrad
and Adam accumulate diagonal statistics
\citep{duchi2011adagrad,kingma2014adam}; and Shampoo-type methods use matrix or
tensor preconditioners \citep{gupta2018shampoo,morwani2024shampoo}.  These
methods differ in how they construct and restrict a positive map from
covectors to parameter directions.

To make the information channel explicit, fix a parameter point \(\theta\),
let \(V=T_\theta\mathcal M\) and \(E=V^*\), and write the per-example loss
differentials as \(a_i=\dd_\theta\ell_i\in E\).  Their span
\[
  W=\operatorname{span}\{a_1,\ldots,a_n\}\subset E
\]
is the first-order visible cotangent space.  Its annihilator
\(\Ann(W)\subset V\) is invisible to every example loss at that point.  A
declared positive visible target \(R\in\SPD(W)\) constrains the restriction of
a full cometric \(P\in\SPD(E)\), but it leaves an entire fiber
\[
  \{P\in\SPD(E):P|_{W\times W}=R\}
\]
undetermined.  The central questions are geometric and matrix-variational:

\begin{quote}
Which positive compressions admit lossless reduction and a canonical shortest
completion, and how do channel motion and structural restrictions change the
completed matrix geometry?
\end{quote}

This question has three levels.  At the first, no linearity or smoothness is
available.  We show that exact reduction of every reference-radial visible
decision problem characterizes metric submetries.  This maximal statement
isolates the precise property required for lossless completion.

At the second level, the spaces are Hadamard manifolds and the information map
is a Riemannian submersion.  If its horizontal distribution integrates to
complete totally geodesic leaves, each leaf is a global isometric information
sheet.  The resulting section is coherent across updates, is the unique
shortest lift, and commutes exactly with proximal maps and gradient flows.

At the third level, the ambient geometry is the SPD cone with the
affine-invariant Riemannian metric (AIRM), and visible information is a
full-column-rank linear channel
\[
  \pi_A(P)=A^\top P A.
\]
This realization is fully explicit.  For reference \(P_0\), set
\(R_0=A^\top P_0A\) and
\(U_0=P_0^{1/2}AR_0^{-1/2}\).  The canonical completion of
\(R\in\SPD^m\) is
\begin{equation}
  \Psi_{P_0,A}(R)
  =P_0^{1/2}
  \left[I+U_0(R_0^{-1/2}RR_0^{-1/2}-I)U_0^\top\right]
  P_0^{1/2}.
  \label{eq:intro-completion}
\end{equation}
It is the unique AIRM-nearest point in the visible fiber, and the full
deformation cost equals the visible cost exactly.

The explicit realization reveals additional structure.  Modulo visible
coordinate changes, active information descriptors \((A,R)\) stratify the
entire SPD cone by \(\rank(P-P_0)\).  In normalized horizontal coordinates,
the exact infinitesimal metric is
\begin{equation}
  \norm{\dot P}_{P,\AI}^2
  =
  \fro{R^{-1/2}\dot R R^{-1/2}}^2
  +2\tr\!\left(K[R+R^{-1}-2I]K^\top\right).
  \label{eq:intro-line-element}
\end{equation}
The first term measures visible geometry change.  The second measures rotation
of the visible subspace.  Reference-valued modes have zero rotation cost and
their admissible rotations form precisely the singular directions of the
descriptor.  Polarizing
\eqref{eq:intro-line-element} gives a closed-form pair metric for arbitrary
coupled perturbation directions.  Its multi-direction matrix is a positive
semidefinite Gram matrix.  This metric cross-term measures the AIRM geometry
of completed motion; it is distinct from the Hessian of the reduced
minimum-distance value.

\paragraph{Contributions.}
The fixed-channel SPD theorem and the moving-channel pair metric are the two
central results.  The abstract, structured, and statistical theorems identify
their mechanism and consequences.
\begin{enumerate}[leftmargin=*,itemsep=2pt]
  \item \emph{Exact SPD reduction and completion.}  We prove that every
  full-column-rank SPD compression is a
  split-Hadamard submetry and derive the closed-form completion
  \eqref{eq:intro-completion}, its horizontal geometry, exact full-to-visible
  decision reduction, and closed-form prior--data shrinkage.  A
  radius-conditioned minimax identity quantifies the remaining invisible
  ambiguity.
  \item \emph{Moving-channel geometry.}  We identify the gauge quotient and
  all rank strata of completed SPD geometry, and derive the closed-form
  mismatch-weighted pair metric whose diagonal is
  \eqref{eq:intro-line-element}.  Its Gram factorization gives exact rank,
  metric-orthogonality and differential-kernel criteria, together with
  moving-reference covariance and the rank-loss boundary.
  \item \emph{General mechanism.}  We characterize metric submetries by
  attained fiber distance and universal radial-visible reduction.  Complete
  horizontal leaves on Hadamard manifolds yield coherent isometric sections,
  exact proximal commutation, and solution-wise gradient-flow lifting.
  \item \emph{Implementability and recovery.}  We characterize diagonal and
  block structured expressivity by
  relative interiors of conic images, with strict separation and facial
  certificates for distinct failure modes.  Deterministic and finite-sample
  bounds then separate visible AIRM error from Grassmann subspace error.
\end{enumerate}

\paragraph{Scope.}
The theory starts after a positive visible target has been declared.  Fisher,
Gauss--Newton, gradient covariance, damping, and user-specified audit targets
provide different constructions of such a target; the paper does not identify
them with one another.  The moving-reference formula is a kinematic identity,
not an optimizer evolution law.  Likewise, exact reduction determines a
canonical full representative of a visible decision but does not prescribe
which visible decision an optimizer should make.

\paragraph{Organization.}
\Cref{sec:related} positions the work.  \Cref{sec:metric-foundations}
establishes the maximal metric theorem, and \cref{sec:split-hadamard} gives its
smooth global form.  \Cref{sec:spd-realization,sec:stratified-geometry}
develop the explicit SPD realization and its rank-stratified pair metric.
\Cref{sec:structured-expressivity} treats structured families, and
\cref{sec:statistical-recovery} develops perturbation and sampling results.
